\documentclass[11pt]{article}
\usepackage[T1]{fontenc}
\usepackage{lmodern}
\usepackage[margin=1in]{geometry}
\usepackage{amsmath,amssymb,amsthm,booktabs}
\usepackage{microtype}
\usepackage[hidelinks]{hyperref}
\hypersetup{pdftitle={A counterexample to multiset rook-Eulerian interlacing},pdfauthor={Xiaojian Zeng}}
\newtheorem{theorem}{Theorem}
\DeclareMathOperator{\asc}{asc}
\newcommand{\W}{\mathcal W}
\title{A counterexample to multiset rook--Eulerian interlacing}
\author{Xiaojian Zeng\thanks{PBC School of Finance, Tsinghua University,
100190. Email: \texttt{xiaojianzeng@foxmail.com}.}}
\date{September 13, 2026}

\begin{document}
\maketitle

\begin{abstract}
Alexandersson, Jal and Quemener conjectured that multiset rook--Eulerian
polynomials on Ferrers boards are real-rooted and that their first-entry
refinements, in decreasing order of the first letter, form an interlacing
sequence. We refute the interlacing clause on the nonrectangular board
$(3,3,3,3,4,4)$ with content $(2,2,1,1)$. Its $60$ admissible words have
explicit refined generating polynomials whose required root order fails
because $-4-\sqrt{13}<-7$. All three refinements and their sum nevertheless
have only real roots. Thus the conjecture as a conjunction is false, while
its separate assertion of universal real-rootedness for the total polynomial
is not decided here. A complete insertion table makes the counting
certificate directly checkable.
\end{abstract}

\section{The assertion and the counterexample}

We use the conventions of Alexandersson--Jal--Quemener
\cite[Definitions 5, 8 and 9; Section 3.3]{AJQ2025}.
A Ferrers board with $n$ rows is described by positive integers
$\lambda_1\leq\cdots\leq\lambda_n$; row $r$ consists of columns
$1,\ldots,\lambda_r$. Let $\alpha=(\alpha_1,\ldots,\alpha_k)$ be a
vector of nonnegative integers with sum $n$. Define
\[
 \W(\lambda,\alpha)=\left\{
 w\in\{1,\ldots,k\}^n:
 \#\{r:w_r=j\}=\alpha_j\ (1\leq j\leq k),\quad
 w_r\leq\lambda_r\ (1\leq r\leq n)
 \right\}.
\]
Each row has one rook, and a column may have several rooks.
For a word $w$, set
\[
 \asc(w)=\#\{r:1\leq r<n,\ w_r<w_{r+1}\}.
\]
These are strict adjacent ascents, with no endpoint sentinel. The total
and first-entry polynomials are
\begin{equation}\label{eq:definitions}
 R(\lambda,\alpha;t)=\sum_{w\in\W(\lambda,\alpha)}t^{\asc(w)},
 \qquad
 R_j(\lambda,\alpha;t)=
 \sum_{\substack{w\in\W(\lambda,\alpha)\\w_1=j}}t^{\asc(w)}.
\end{equation}
In particular, $R=\sum_{j=1}^{\lambda_1}R_j$, with empty sums equal to zero.

For polynomials $f,g$ with positive leading coefficients and real
nonpositive roots, write their roots, with multiplicity, in decreasing
order as $a_1,a_2,\ldots$ and $b_1,b_2,\ldots$, respectively.
The convention $f\preceq g$ means
\begin{equation}\label{eq:interlace}
 \cdots\leq a_3\leq b_3\leq a_2\leq b_2\leq a_1\leq b_1\leq0,
 \qquad \deg g-\deg f\in\{0,1\}.
\end{equation}
An interlacing sequence $(f_1,\ldots,f_m)$ must satisfy
$f_i\preceq f_j$ for \emph{every} $i<j$.
Conjecture 28 of \cite{AJQ2025}, in its arXiv v1 formulation, asserts
real-rootedness of $R(\lambda,\alpha;t)$ for every Ferrers board and content,
together with interlacing of
\begin{equation}\label{eq:conjectured}
 (R_{\lambda_1},R_{\lambda_1-1},\ldots,R_2,R_1).
\end{equation}

\begin{theorem}\label{thm:witness}
For $\lambda=(3,3,3,3,4,4)$ and $\alpha=(2,2,1,1)$, there are exactly
$60$ words in $\W(\lambda,\alpha)$, and
\begin{align}\label{eq:refinements}
 R_3(t)&=t(t^2+8t+3),\notag\\
 R_2(t)&=t(8t^2+15t+1),\\
 R_1(t)&=t^2(t+7)(2t+1).\notag
\end{align}
The sequence $(R_3,R_2,R_1)$ is not interlacing. Each of these three
polynomials and the total
\begin{equation}\label{eq:total}
 R(t)=2t(t^3+12t^2+15t+2)
\end{equation}
is real-rooted.
\end{theorem}

The board is non-skew and nonrectangular, and every content coordinate is
positive. The theorem therefore refutes the interlacing clause of the
stated conjecture without a degeneracy involving zero refinements. It
does not refute the separately quantified total-real-rootedness clause.

The relevant preceding results concern two different specializations.
In the nonattacking, distinct-letter setting,
\cite[Theorem 12]{AJQ2025} proves refinement interlacing and total
real-rootedness. For rectangles, the context is Simion's theory of
multiset composition and Eulerian polynomials \cite{Simion1984}; the
identification of rectangular total real-rootedness with Section 2 of that paper
is recorded in \cite[Section 3.3]{AJQ2025}.
Ma and Pan \cite[Theorem 1.11]{MaPan2023} prove first-entry interlacing
for unrestricted multiset words on at least three distinct letters, using
descents. Complementing the alphabet exchanges ascents with descents
and reverses the first-entry order, so their result supplies the
corresponding rectangular comparison relevant here. More recently,
Yan \cite[Theorems 1.4 and 1.5]{Yan2026} proves real-rootedness for
unrestricted colored multiset totals and interlacing when one multiplicity
increases. Those statements concern different word sets and comparisons
from \eqref{eq:conjectured} with row bounds.
The present contribution is the explicit failure of the proposed
first-entry extension when repeated letters and nonrectangular row bounds
occur together. We found no equivalent counterexample in the primary
sources inspected through September 13, 2026; this is a bounded literature
comparison, not an exhaustive priority claim.

\section{Counting and root order}

\begin{proof}[Proof of Theorem~\ref{thm:witness}]
The unique letter $4$ can occupy only position $5$ or $6$. Deleting it
leaves a permutation $v=v_1\cdots v_5$ of $11223$. Conversely, each such
base gives exactly the two admissible words
\[
 v_1v_2v_3v_4\,4\,v_5,
 \qquad
 v_1v_2v_3v_4v_5\,4.
\]
Deletion is inverse to insertion, so no word is missed or counted twice.
There are $5\binom42=30$ bases: choose the position of $3$, then the
two positions of $1$, and fill the remaining positions with $2$.
Thus there are exactly $60$ admissible words.

Write $a_5(v)$ and $a_6(v)$ for the ascent counts after inserting $4$
in position $5$ and $6$, respectively. Since every letter of $v$ is
less than $4$,
\begin{equation}\label{eq:insertion}
 a_5(v)=1+\sum_{r=1}^3\mathbf{1}_{\{v_r<v_{r+1}\}},
 \qquad
 a_6(v)=1+\sum_{r=1}^4\mathbf{1}_{\{v_r<v_{r+1}\}}.
\end{equation}
The complete base table is in Appendix~\ref{app:certificate}.
Grouping its entries by the first letter gives the following two
contributions:
\begin{center}
\begin{tabular}{ccc}
\toprule
First letter & $4$ in position $5$ & $4$ in position $6$\\
\midrule
$1$ & $5t^2+7t^3$ & $2t^2+8t^3+2t^4$\\
$2$ & $t+9t^2+2t^3$ & $6t^2+6t^3$\\
$3$ & $2t+4t^2$ & $t+4t^2+t^3$\\
\bottomrule
\end{tabular}
\end{center}
Adding the two columns and factoring proves \eqref{eq:refinements}.

In particular, the roots of $R_3$ and $R_1$, with multiplicity, are
\[
 \{-4-\sqrt{13},-4+\sqrt{13},0\},
 \qquad
 \{-7,-\tfrac12,0,0\}.
\]
Since $\deg R_3=3$ and $\deg R_1=4$, the necessary chain for
$R_3\preceq R_1$ starts with
$\min\operatorname{root}(R_1)\leq\min\operatorname{root}(R_3)$.
This fails:
\begin{equation}\label{eq:obstruction}
 -4-\sqrt{13}<-7,
 \qquad\text{because }13>3^2.
\end{equation}
The pair $(R_3,R_1)$ is required by the all-pairs definition of an
interlacing sequence. Hence \eqref{eq:conjectured} fails.
There is no cancellation of roots at zero in this comparison.

The displayed roots show that $R_1$ and $R_3$ are real-rooted.
The two nonzero roots of $R_2$ are
$(-15\pm\sqrt{193})/16$, which are real and negative because
$0<193<15^2$. Summing \eqref{eq:refinements} gives \eqref{eq:total}.
For $q(t)=t^3+12t^2+15t+2$, exact integer evaluation gives
\[
\begin{array}{c|rrrrr}
x&-11&-10&-2&-1&0\\\hline
q(x)&-42&52&12&-2&2
\end{array}
\]
There is a root in each of $(-11,-10)$, $(-2,-1)$ and $(-1,0)$ by
the intermediate value theorem. These three distinct roots exhaust the
degree of $q$. The additional root of $R=2tq$ is zero, proving the last
assertion.
\end{proof}

The root certificate uses exact order comparisons and polynomial
evaluations. The same type of certification is used for an algebraically
defined bound by Zeng, Liu and Ratnavelu
\cite{ZengLiuRatnavelu2026}, Section 2, equations (17)--(18), where a
rational evaluation of a cubic and monotonicity replace a decimal root
estimate. Here the obstruction reduces to \eqref{eq:obstruction}, and
integer sign changes certify the total. No numerical root approximation
is needed.

The example establishes neither a minimal counterexample size nor a
classification of boards admitting interlacing. In particular, the
universal question whether every nonzero total
$R(\lambda,\alpha;t)$ has only real roots remains unresolved by this
argument.

\appendix
\section{Complete insertion certificate}\label{app:certificate}

The $30$ bases and their two ascent counts in \eqref{eq:insertion} are
listed below. There are $12$ bases beginning with $1$, $12$ beginning
with $2$, and $6$ beginning with $3$. The two halves of the table together
list every permutation of $11223$ exactly once.

\begin{center}
\begin{tabular}{crr@{\qquad\qquad}crr}
\toprule
$v$ & $a_5(v)$ & $a_6(v)$ & $v$ & $a_5(v)$ & $a_6(v)$\\
\midrule
11223 & 2 & 3 & 21231 & 3 & 3\\
11232 & 3 & 3 & 21312 & 2 & 3\\
11322 & 2 & 2 & 21321 & 2 & 2\\
12123 & 3 & 4 & 22113 & 1 & 2\\
12132 & 3 & 3 & 22131 & 2 & 2\\
12213 & 2 & 3 & 22311 & 2 & 2\\
12231 & 3 & 3 & 23112 & 2 & 3\\
12312 & 3 & 4 & 23121 & 3 & 3\\
12321 & 3 & 3 & 23211 & 2 & 2\\
13122 & 3 & 3 & 31122 & 2 & 2\\
13212 & 2 & 3 & 31212 & 2 & 3\\
13221 & 2 & 2 & 31221 & 2 & 2\\
21123 & 2 & 3 & 32112 & 1 & 2\\
21132 & 2 & 2 & 32121 & 2 & 2\\
21213 & 2 & 3 & 32211 & 1 & 1\\
\bottomrule
\end{tabular}
\end{center}

As a separate check, one may visit all $4^6=4096$ words on
$\{1,2,3,4\}$, retain content $(2,2,1,1)$, and impose the row bounds.
Exactly $180$ words pass the content test and $60$ pass both tests.
The following complete Python procedure computes coefficient lists in
increasing degree order; it requires only the standard library.
\begin{quote}\small
\begin{verbatim}
from itertools import product
lam = (3, 3, 3, 3, 4, 4)
alpha = (2, 2, 1, 1)
c = {j: [0] * 6 for j in (1, 2, 3)}
for w in product(range(1, 5), repeat=6):
    if tuple(w.count(j) for j in range(1, 5)) != alpha:
        continue
    if any(x > b for x, b in zip(w, lam)):
        continue
    a = sum(w[r] < w[r + 1] for r in range(5))
    c[w[0]][a] += 1
assert c == {1: [0, 0, 7, 15, 2, 0],
             2: [0, 1, 15, 8, 0, 0],
             3: [0, 3, 8, 1, 0, 0]}
\end{verbatim}
\end{quote}
The enumeration was also implemented by choosing the base positions
and by visiting all $6!$ labeled orders, erasing the labels on equal
letters, and filtering the row bounds. The resulting word sets and ascent
counts agree. These computations check the finite table; the insertion
bijection and root comparison give the mathematical argument.

\section*{Computational and AI assistance}
AI tools assisted with witness exploration, programming, proof drafting,
source searches and manuscript preparation. The finite count was
recomputed by separately implemented enumeration methods, including a
Cartesian-product check that imports no discovery code. All decisive
counting and sign computations use exact integer or rational arithmetic.
The proof above is explicit and does not depend on a numerical fit,
a root plot, or a formal proof-assistant verification.

\bibliographystyle{plain}
\raggedright
\bibliography{references}
\end{document}
